\section{Easter Prayer Through the Day}

\subsection{From Easter Sunday through Pentecost}

The Directory 196 says the Regina Coeli replaces the Angelus during Easter Time. Under the current Roman calendar, Easter Time comprises the fifty days from Easter Sunday through Pentecost Sunday inclusive. This present boundary differs from the 1925 missal's historical direction, which ran from noon on Holy Saturday until noon on the Saturday before Trinity Sunday. Follow the current calendar and competent local books rather than silently carrying that older schedule forward.

The ordinary daily rhythm remains early morning, midday, or toward evening. These are recurring thresholds, not universally fixed clock minutes. A sustainable practice may connect the prayer with the beginning of work, a midday pause, and the close of labor. The devotion is recommended rather than imposed on every Christian; missing a time is not a violation of a universal precept.

\subsection{Ordinary, solemn, and liturgical forms}

One person may say both the versicle and response, while a group may alternate them. A community may sing the antiphon. The Directory also proposes that a more solemn popular celebration may include proclamation of a Resurrection Gospel. Such enrichment should direct attention to the Paschal Mystery rather than turn the exercise into a substitute Mass or an improvised liturgical hour.

The 1925 witness says the prayer is recited standing. Standing remains an intelligible Easter custom and bodily sign of resurrection joy, but the current indulgence grant does not make posture a condition. Likewise, a bell may summon a community to prayer without being necessary to the prayer's validity or current indulgenced use.

When the Regina caeli is prayed as part of the Liturgy of the Hours, the governing liturgical book controls its text, translation, placement, and manner. When its expanded form replaces the Angelus, it is a popular exercise. The shared antiphon does not erase that distinction.

\subsection{Current indulgence discipline}

As checked on 21 July 2026, the Apostolic Penitentiary's fourth \textit{Enchiridion Indulgentiarum} (1999), concession 17 \S2, 2\textsuperscript{o}, grants a \emph{partial} indulgence to a member of the faithful who piously recites the Angelus with its appointed exchanges and collect early in the morning, at midday, or toward evening; during Easter Time it names the Regina caeli with its customary prayer instead.

The Latin alternatives mean that one pious recitation at one of the named times is the prescribed work. All three times form the recommended traditional rhythm, but they are not a set that must be completed before any indulgence can be gained. The grant remains subject to the \textit{Enchiridion}'s general norms: the faithful must be baptized, not excommunicated, and in the state of grace at least when the work is completed; have at least the general intention of gaining the indulgence; perform the work duly; and, for a partial indulgence, act with at least a contrite heart (norms 4 and 17). A vernacular prayer used for the indulgence requires approval by competent ecclesiastical authority (norm 22). These are not the additional conditions proper to a plenary indulgence.

The current grant states no requirement to hear a bell, stand, pray in a church, or add the Glory Be. It does not create an obligation to recite the devotion.

Lefebvre's 1925 rubric says that the same indulgences could be gained as for the Angelus and directs standing. Later historical collections describe quantified partial indulgences and periodic plenary grants under conditions then applicable. Paul VI's \textit{Indulgentiarum doctrina} (1967) abolished the former measurement of partial indulgences in days or years and required a revised collection. Those superseded measurements and conditions must not be represented as the present grant.

An indulgence is an exercise of the Church's ministry within the life of conversion. It is not forgiveness purchased by recitation, a calculation of time in purgatory, or permission to neglect sacramental reconciliation, restitution, amendment of life, or charity.
